\section{A complex-surface lift that retains the Euclidean metric}

The algebraic complexification above has complex dimension three. There is a different, exact construction on a complex surface: add one auxiliary circle to physical three-space. This is a static gauge-theoretic lift, not an identification of that circle with plasma evolution time.

Let $\Omega$ be a smooth solid torus, let $B$ be a smooth real field on a neighborhood of its closure, and set
\[
\alpha=B^\flat,\qquad \beta=*_{3}\alpha,\qquad
M=\Omega\times(\mathbb R/\ell\mathbb Z).
\]
Give $M$ the product Euclidean metric and orientation $dx\wedge dy\wedge dz\wedge dt$. Define the real two-form
\[
\mathcal F=\beta-\alpha\wedge dt.
\]

\paragraph{Proposition (exact Abelian Hermitian--Yang--Mills lift).}
The field $B$ is vacuum if and only if $d\mathcal F=0$. The form $\mathcal F$ is anti-self-dual for every $B$. Consequently a vacuum field gives a closed primitive form of type $(1,1)$ for the complex coordinates
\[
u=x+iy,\qquad w=\exp(2\pi(z+it)/\ell).
\]
If $\Omega$ avoids the vertical axis, $M$ is a complex domain in $(\mathbb C^*)^2$. Its metric and Kähler form are
\[
g_4=|du|^2+\frac{\ell^2}{4\pi^2|w|^2}|dw|^2,
\qquad
\omega=\frac i2du\wedge d\bar u+
\frac{i\ell^2}{8\pi^2|w|^2}dw\wedge d\bar w.
\]
In particular, the physical metric is not replaced by the standard flat Hermitian metric in the algebraic coordinates $(u,w)$.

\paragraph{Proof.}
The components of $d\mathcal F=d\beta-d\alpha\wedge dt$ with and without $dt$ vanish separately. Their vanishing says respectively $\operatorname{div}B=0$ and $\operatorname{curl}B=0$. The product Hodge star satisfies
\[
*_{4}\beta=\alpha\wedge dt,\qquad
*_{4}(\alpha\wedge dt)=\beta.
\]
Thus $*_{4}\mathcal F=-\mathcal F$. On this Kähler surface, real anti-self-dual two-forms are precisely the real primitive $(1,1)$ forms. This last assertion can also be verified in the orthonormal coordinates $(x,y,z,t)$; no compactness or algebraic compactification is used. The metric formula follows by differentiating $w$. $\square$

\paragraph{Integral periods and the holomorphic line bundle.}
Let $\gamma$ generate the longitudinal first homology of $\Omega$, and write $\Gamma=\int_\gamma\alpha$. Since $H^2(\Omega;\mathbb R)=0$, the only possible nonzero period of $\mathcal F$ on $M$ is
\[
\frac{1}{2\pi}\int_{\gamma\times S^1_\ell}\mathcal F
=-\frac{\Gamma\ell}{2\pi}.
\]
For the normalized seed, $\Gamma=2\pi$ and $\ell=1$ make this period integral. More generally one may choose the auxiliary period or an overall mathematical curvature normalization. This is not a physical flux-quantization assertion.

Choose a global one-form $a$ on $\Omega$ with $da=\beta$. Locally write $\alpha=d\Phi_i$ and set $A_i=a-\Phi_i\,dt$. Then $dA_i=\mathcal F$. A potential branch change by $2\pi n$ changes $A_i$ by $-2\pi n\,dt$, which has the unitary transition $e^{2\pi int}$ when $\ell=1$. With convention $\nabla=d-iA$, the curvature is $-i\mathcal F$ and $c_1=[\mathcal F/(2\pi)]$. The type $(1,1)$ condition makes the $(0,1)$ part of this connection integrable, hence defines a holomorphic line bundle; the primitive condition is the Abelian Hermitian--Yang--Mills equation with zero Einstein constant. The domain need not be assumed Stein for this explicit construction.

The logarithmic differential on the threefold is algebraic, but this line bundle on the surface need not be algebraic. The algebraic Picard group of $(\mathbb C^*)^2$ is zero: its coordinate ring is a Laurent polynomial unique factorization domain and is regular. A bundle with the nonzero Chern period above therefore cannot be the restriction of an algebraic line bundle on that ambient algebraic torus. This is a genuine distinction between the two constructions, not an inconsistency.

\paragraph{A positive auxiliary energy and its exact scope.}
For a vacuum field,
\[
\int_M|\mathcal F|^2=2\ell\int_\Omega|B|^2.
\]
If $a_1$ is a real one-form whose tangential trace on $\partial M$ vanishes, then $d*\mathcal F=0$ and integration by parts give
\[
\int_M|\mathcal F+da_1|^2-\int_M|\mathcal F|^2
=\int_M|da_1|^2\ge0.
\]
Thus the lift minimizes the Abelian Yang--Mills energy against these fixed-boundary connection variations. A moving plasma boundary, plasma pressure, distribution functions, collisions, and source work are absent from this identity. It is not a plasma Lyapunov theorem.

\section{A determinant quantity for vacuum geometry}

Here is a separate elementary determinant construction that retains the physical Euclidean Hessian. It gives an exact positive identity and useful geometric control, although not a dynamical plasma estimate.

\paragraph{Proposition (a subharmonic log-determinant).}
Let $B$ be a smooth vacuum field on a neighborhood of the closure of a bounded smooth domain. Set $H=\nabla B$, choose $\lambda>0$, and suppose $\|H\|_{\mathrm{op}}\le\delta\lambda$ on the boundary, with $0<\delta<1$. Then throughout the domain
\[
A=I+\lambda^{-1}H>0,\qquad
Q=-\log\det A\ge0,
\]
and
\[
\Delta Q=\sum_k\operatorname{tr}
\big(A^{-1}(\partial_k A)A^{-1}(\partial_k A)\big)\ge0.
\]
Moreover,
\[
\frac{|H|_F^2}{2\lambda^2(1+\delta)^2}
\le Q\le
\frac{|H|_F^2}{2\lambda^2(1-\delta)^2},
\qquad
\Delta Q\ge\frac{|\nabla H|_F^2}{\lambda^2(1+\delta)^2}.
\]

\paragraph{Proof.}
Vacuum implies $H^T=H$, $\operatorname{tr}H=0$, and $\Delta H=0$. For every fixed real vector $\xi$, the scalar $\xi^TH\xi$ is harmonic. The maximum principle therefore extends the boundary spectral bounds $-\delta\lambda I\le H\le\delta\lambda I$ throughout the domain. Differentiating $-\log\det A$, and using $\Delta A=0$, gives the displayed Laplacian identity. Each summand is the squared Frobenius norm of $A^{-1/2}(\partial_kA)A^{-1/2}$. For the comparison bounds, let $h_j$ be the eigenvalues of $H/\lambda$. Since $\sum_jh_j=0$,
\[
Q=\sum_j\big(h_j-\log(1+h_j)\big).
\]
The scalar function in parentheses has value and first derivative zero at the origin, and second derivative $(1+h)^{-2}$. Taylor's integral formula on $[-\delta,\delta]$ gives the two-sided bounds. The lower bound for $\Delta Q$ follows by the same spectral bounds on $A^{-1/2}$. $\square$

Locally, $B=\nabla\Phi$ and $A$ is the Hessian of $|x|^2/2+\Phi/\lambda$. Different branches of the angular potential differ by constants, so $A$ and $Q$ remain globally well-defined. This is a real Monge--Ampère-type quantity with a fully checked sign. The maximum principle controls $Q$ from its boundary data. It does not require, or assert, a globally single-valued convex potential on the solid torus.

The limitations are explicit. A general MHD or kinetic perturbation has current; then $\nabla B$ need not be symmetric and $\Delta\nabla B$ need not vanish. The displayed identity does not close a plasma evolution estimate. A successful auxiliary-PDE argument would have to include those terms and the physical boundary terms with controlled signs.

\paragraph{A tempting simplification is impossible.}
A nonzero smooth vacuum field tangent to the boundary of a bounded smooth domain cannot have constant strength on the entire boundary. To prove this, put $q=|B|^2$. Then $\Delta q=2|\nabla B|^2\ge0$. If $q=c>0$ on the boundary and is nonconstant inside, the maximum principle and Hopf lemma give $\partial_nq>0$ at every boundary point. Tangency and curl-freeness instead give
\[
\frac12\partial_nq=-\mathrm{II}(B,B),
\]
where $\mathrm{II}(v,w)=\langle D_vn,w\rangle$. At a point maximizing distance from a fixed interior point, comparison with the enclosing sphere gives $\mathrm{II}\ge R^{-1}g$ on the tangent plane. There $B\ne0$, so the last normal derivative is negative, a contradiction. If $q$ is constant throughout, the Bochner identity forces $\nabla B=0$; a nonzero constant vector cannot be tangent to the boundary of a bounded domain, as integrating $\operatorname{div}((B\cdot x)B)=|B|^2$ shows. The case $c=0$ forces $B=0$ by the maximum principle. This excludes that particular attempt to remove field-strength variation; it is not an instability theorem.
